\section{Discussion}

\subsection{The connected laboratory is maintained through boundary work}

The findings show that maintenance in connected laboratories is organized around boundaries among representations, actors, and forms of evidence. Participants must translate between vendor and open-source vocabularies, source files and runtime state, modeled and physical resources, logs and scientific outcomes, public discussion and restricted support. Breakdowns make these dependencies visible, and collaborative work temporarily assembles a representation adequate for action.

This account extends an infrastructure perspective in two ways. First, the relevant infrastructure includes epistemic boundaries. Participants negotiate what evidence supports a claim, which validation stage has been reached, and whether a local result can travel to another device or workflow. Second, the infrastructure includes material irreversibility. The coordinated object after an incident is not simply a software state. It is a partially transformed physical process whose history cannot be undone by restoring code.

\subsection{Articulation work is also requirements production}

Participants do not merely supply missing details requested by an expert. Through discussion they establish the requirement set itself. Scientific timing, storage, contamination limits, physical feedback, licensing, and state synchronization enter because community members recognize dependencies omitted from the initiating formulation. The forum functions as a distributed requirements-elicitation system.

This process produces value even when the local problem remains unresolved. The scheduling discussion never reached a fully specified instance, although it generated a clearer representation of the fields needed for comparable scheduling. The AI discussion did not establish a production method generator, although it produced a staged account of context and validation. Public technical communities should therefore be evaluated on diagnostic sufficiency and requirement discovery alongside answer acceptance.

\subsection{Repair needs a record of uncertainty}

Repair scholarship emphasizes restoring workable sociotechnical order. The laboratory cases show that repair also requires representing uncertainty about the physical world. The final confirmed command can differ from the final physical action. A sensor can provide an observation without establishing transferred volume. Manual intervention can restore operation and invalidate a sample. A recovery process therefore needs explicit unknown and uncertain state, not only a success or failure flag.

The absence of final disposition across the observability cases is analytically important. Support systems commonly optimize for error identification and resumption. Laboratory repair additionally requires a judgment about material and scientific result. This judgment can involve expertise, risk tolerance, sample value, and local policy. Systems that expose the uncertainty and provenance of that decision can support collaborative repair without pretending to automate it completely.

\subsection{Knowledge migration creates a public maintenance deficit}

The movement from public forum to private channel should not be understood only as lost data. It reflects legitimate constraints: commercial documentation, confidential protocols, remote debugging, restricted software, or the need for synchronous collaboration. The problem arises when public discovery leads to private resolution without a bounded public return. Future practitioners can find the original question and not the current answer.

A public maintenance model should therefore support return paths. After private work, participants should be able to publish a redacted resolution, applicability statement, evidence grade, limitation, and maintained artifact. The source discussion can retain credit and context. A canonical index can expose the current state without requiring disclosure of proprietary details. The challenge is organizational and representational, not solely motivational.

\subsection{Metrics can support reflection without ranking communities}

The bounded metrics helped identify where evidence was incomplete and which relationships deserved prospective study. Their value lies in making definitions contestable. Integration accessibility separates documentation, protocol, licensing, testing, examples, and maintenance. Observability separates command, observation, state, intervention, recovery, and disposition. Test--claim alignment separates evidence stages.

These metrics should not become performance rankings for vendors or community members. They are instruments for reflecting on whether a case record contains the information required for a specified judgment. Their small denominators and purposive sampling make this boundary explicit. Similar instruments can support other technical communities when the local units and evidence rules are defined with participants.
